\section*{Capítulo \ref{ch:b2:frontier-orbitals} --- Orbitais de fronteira e reatividade}

\begin{solution}{exo:b2:frontier-orbitals:1}
Eteno: \omterm{def:b2:frontier-orbitals:frontier}{HOMO} $\alpha + \beta$, \omterm{def:b2:frontier-orbitals:frontier}{LUMO} $\alpha - \beta$. Ânion alila (quatro elétrons): \omterm{def:b2:frontier-orbitals:frontier}{HOMO} $\alpha$ (o
nível não ligante), \omterm{def:b2:frontier-orbitals:frontier}{LUMO} $\alpha - 1.414\beta$. Butadieno: \omterm{def:b2:frontier-orbitals:frontier}{HOMO} $\alpha + 0.618\beta$, \omterm{def:b2:frontier-orbitals:frontier}{LUMO} $\alpha -
0.618\beta$.
\end{solution}

\begin{solution}{exo:b2:frontier-orbitals:2}
Duros: \ce{OH-}, \ce{H+}, \ce{F-}. Moles: \ce{I-}, \ce{RS-}, o carbono do iodometano.
\end{solution}

\begin{solution}{exo:b2:frontier-orbitals:3}
O butadieno (confôrmero s-cis acessível) e o ciclopentadieno (travado em s-cis)
reagem; o penta-1,4-dieno não é conjugado e não reage. O \omterm{def:b2:frontier-orbitals:diels-alder}{dieno} (2\textit{Z},4\textit{Z}) só
consegue atingir a forma s-cis empurrando seus dois grupos metila internos um
contra o outro: ele reage muito mal.
\end{solution}

\begin{solution}{exo:b2:frontier-orbitals:4}
Ciclo-hex-3-eno-1-carbonitrila: um anel de seis membros, com a ligação dupla
entre os antigos C2 e C3 do \omterm{def:b2:frontier-orbitals:diels-alder}{dieno} e o grupo \ce{C#N} em um dos dois carbonos
do anel formados a partir do \omterm{def:b2:frontier-orbitals:diels-alder}{dienófilo}.
\end{solution}

\begin{solution}{exo:b2:frontier-orbitals:5}
$\Delta\varepsilon = \qty{8}{eV}$: perturbativa $2\beta^2/\Delta\varepsilon = \qty{0.25}{eV}$. Exata: nível
inferior $-6 - \sqrt{16 + 1} = \qty{-10.123}{eV}$, ganho $2 \times 0.123 = \qty{0.246}{eV}$.
\end{solution}

\begin{solution}{exo:b2:frontier-orbitals:6}
O iodometano é um \omterm{def:b2:frontier-orbitals:hard-soft}{eletrófilo mole}: \omterm{def:b2:frontier-orbitals:control}{controle orbital}, ataque pelo átomo de
maior coeficiente no \omterm{def:b2:frontier-orbitals:frontier}{HOMO}, o carbono. Um carbocátion livre é um \omterm{def:b2:frontier-orbitals:hard-soft}{eletrófilo
duro}: \omterm{def:b2:frontier-orbitals:control}{controle de carga}, ataque em parte pelo nitrogênio, que tem mais carga
negativa.
\end{solution}

\begin{solution}{exo:b2:frontier-orbitals:7}
O C4 do \omterm{def:b2:frontier-orbitals:diels-alder}{dieno} (maior coeficiente no \omterm{def:b2:frontier-orbitals:frontier}{HOMO}) se liga ao carbono $\beta$ do propenal
(maior coeficiente no \omterm{def:b2:frontier-orbitals:frontier}{LUMO}): 2-metoxiciclo-hex-3-eno-1-carbaldeído, com os dois
substituintes em carbonos vizinhos (produto ``orto'').
\end{solution}

\begin{solution}{exo:b2:frontier-orbitals:8}
Um esqueleto biciclo[2.2.1]hepteno (norborneno) com o grupo éster em um carbono
do antigo \omterm{def:b2:frontier-orbitals:diels-alder}{dienófilo}, voltado para a nova ligação \ce{C=C}, sob o anel de seis
membros (endo), do lado oposto à ponte \ce{CH2}.
\end{solution}

\begin{solution}{exo:b2:frontier-orbitals:9}
Os dois grupos carbonila abaixam o \omterm{def:b2:frontier-orbitals:frontier}{LUMO} do \omterm{def:b2:frontier-orbitals:diels-alder}{dienófilo}: sua diferença com o \omterm{def:b2:frontier-orbitals:frontier}{HOMO}
do ciclopentadieno é muito menor que a do eteno, a estabilização
$\beta^2/\Delta\varepsilon$ do estado de transição é maior, e a barreira, menor.
\end{solution}

\begin{solution}{exo:b2:frontier-orbitals:10}
A reação é estereoespecífica: o maleato dá o aduto com os dois grupos éster cis
(endo,endo como produto majoritário; ele é aquiral, um composto meso); o
fumarato dá o aduto trans, com um éster endo e um exo, formado como um par de
enantiômeros (racêmico).
\end{solution}

\begin{solution}{exo:b2:frontier-orbitals:11}
Duas ligações $\pi$ se tornam duas ligações $\sigma$ mais fortes: $\Delta_r H^\circ < 0$.
Duas moléculas se tornam uma:
$\Delta_r S^\circ < 0$. $\Delta_r G^\circ = \Delta_r H^\circ - T\Delta_r S^\circ$ cresce com $T$ e se torna
positivo acima de $T_i = \Delta_r H^\circ/\Delta_r S^\circ$: a alta temperatura, o aduto se
decompõe de volta (retro-Diels--Alder).
\end{solution}

\begin{solution}{exo:b2:frontier-orbitals:12}
Dois orbitais de pares não ligantes preenchidos lado a lado formam uma
interação a quatro elétrons, que é repulsiva. A molécula gira em torno da
ligação \ce{N-N} até que os pares não ligantes fiquem aproximadamente
perpendiculares (uma conformação gauche), onde sua sobreposição, e a repulsão,
são pequenas.
\end{solution}

\begin{solution}{pb:b2:frontier-orbitals:1}
\textbf{1.} Um anel de cinco membros com duas ligações duplas conjugadas e um
\ce{CH2}: o anel mantém a unidade \omterm{def:b2:frontier-orbitals:diels-alder}{dieno} na conformação s-cis.
\textbf{2.} O \omterm{def:b2:frontier-orbitals:diels-alder}{dienófilo}, por uma de suas ligações duplas.
\textbf{3.} Um esqueleto norborneno fundido a um anel ciclopenteno; as duas
novas ligações $\sigma$ unem as extremidades do \omterm{def:b2:frontier-orbitals:diels-alder}{dieno} aos dois carbonos da
ligação dupla do \omterm{def:b2:frontier-orbitals:diels-alder}{dienófilo}.
\textbf{4.} O \omterm{def:b2:frontier-orbitals:endo}{aduto endo} (\omterm{def:b2:frontier-orbitals:endo}{regra endo}): o anel restante do \omterm{def:b2:frontier-orbitals:diels-alder}{dienófilo} fica sob o
\omterm{def:b2:frontier-orbitals:diels-alder}{dieno} no estado de transição.
\textbf{5.} Sim: ela é concertada, as duas ligações se formam na mesma face de
cada parceiro, e a geometria do \omterm{def:b2:frontier-orbitals:diels-alder}{dienófilo} é conservada.
\textbf{6.} O ciclopentadieno é ao mesmo tempo um bom \omterm{def:b2:frontier-orbitals:diels-alder}{dieno} (travado em s-cis) e
um \omterm{def:b2:frontier-orbitals:diels-alder}{dienófilo}, e a reação é concertada, sem intermediário carregado a
estabilizar.
\textbf{7.} Negativo: duas ligações $\pi$ são substituídas por duas ligações $\sigma$.
\textbf{8.} Negativo: duas moléculas dão uma.
\textbf{9.} $\Delta_r G^\circ = \Delta_r H^\circ - T\Delta_r S^\circ$ com os dois termos negativos: negativo a
$T$ baixa, nulo em $T_i = \Delta_r H^\circ/\Delta_r S^\circ$, positivo acima.
\textbf{10.} À temperatura ambiente, $T < T_i$: a dimerização é favorecida, embora
lenta. Perto do ponto de ebulição do dímero, o equilíbrio favorece o monômero
e é atingido depressa.
\textbf{11.} O monômero deixa a mistura quente como vapor: retirar um produto
mantém o equilíbrio se deslocando para o monômero (Le Chatelier).
\textbf{12.} À temperatura ambiente, ele torna a dimerizar; o frio desacelera
essa reação.
\textbf{13.} $0.62 - (-0.62) = 1.24|\beta|$.
\textbf{14.} $0.62 - (-0.35) = 0.97|\beta|$.
\textbf{15.} Ciclopentadieno com anidrido maleico: a diferença menor dá uma
estabilização maior do estado de transição, uma reação mais rápida.
\textbf{16.} Eles estão conjugados com a ligação \ce{C=C} e são retiradores de
elétrons: como no propenal, eles abaixam o nível $\pi^*$.
\textbf{17.} Seu \omterm{def:b2:frontier-orbitals:frontier}{LUMO}, que tem lobos nos carbonos das carbonilas diante dos
carbonos centrais do \omterm{def:b2:frontier-orbitals:frontier}{HOMO} do \omterm{def:b2:frontier-orbitals:diels-alder}{dieno}.
\textbf{18.} Esqueleto norborneno; o anel do anidrido sob a nova ligação dupla,
do lado oposto à ponte \ce{CH2}.
\textbf{19.} Os dois carbonos cabeça de ponte e os dois carbonos que levam o
anidrido: quatro estereocentros.
\textbf{20.} Não: um plano de simetria passa pela ponte \ce{CH2} e entre as duas
metades da molécula; é um composto meso.
\textbf{21.} Eles estavam cis no \omterm{def:b2:frontier-orbitals:diels-alder}{dienófilo}, e a adição concertada os mantém cis.
\textbf{22.} O mesmo esqueleto, com o anel do anidrido do lado da ponte \ce{CH2}.
\textbf{23.} Não: passar de endo a exo exigiria romper e refazer ligações
\ce{C-C}, pela reação inversa, que exige aquecimento.
\textbf{24.} A água abre o anidrido: o diácido carboxílico cis, com os dois
grupos \ce{COOH} endo.
\textbf{25.} O \omterm{def:b2:frontier-orbitals:endo}{aduto endo} tem $\boldsymbol{4}$ estereocentros.
\end{solution}
